% Part: first-order-logic
% Chapter: natural-deduction
% Section: proving-things-quant

\documentclass[../../../include/open-logic-section]{subfiles}

\begin{document}

\olfileid{fol}{ntd}{prq}

\olsection{\usetoken{P}{derivation} با سورها}

\begin{ex}
هنگام کار با سورها باید مراقب باشیم شرط متغیر ویژه نقض نشود، و گاهی
این امر ایجاب می‌کند ترتیب انجام برخی استنتاج‌ها را تغییر دهیم. به‌طور
کلی، بهتر است نخست به قواعدی بپردازیم که مشمول شرط متغیر ویژه‌اند
(آن‌ها در برهان تکمیل‌شده پایین‌تر خواهند بود).

ببینیم چگونه می‌توانیم !!a{derivation} از !!{formula}
$\lexists[x][\lnot !A(x)] \lif \lnot \lforall[x][!A(x)]$ به دست دهیم.
طبق معمول با نوشتن عبارت زیر آغاز می‌کنیم:
\begin{prooftree}
\AxiomC{}
\UnaryInfC{$\lexists[x][\lnot !A(x)]\lif \lnot \lforall[x][!A(x)]$}
\end{prooftree}
ابتدا می‌نویسیم برای توجیه آن گام آخر با استفاده از قاعدهٔ
\Intro{\lif} چه چیزی لازم است.
\begin{prooftree}
\AxiomC{$\Discharge{\lexists[x][\lnot !A(x)]}{1}$}
\DeduceC{$\lnot \lforall[x][!A(x)]$}
\DischargeRule{\Intro{\lif}}{1}
\UnaryInfC{$\lexists[x][\lnot !A(x)]\lif \lnot \lforall[x][!A(x)]$}
\end{prooftree}
چون قاعدهٔ آشکاری برای اعمال بر $\lnot \lforall[x][!A(x)]$ وجود
ندارد، !!{derivation} را چنان ترتیب می‌دهیم که بتوانیم از قاعدهٔ
\Elim{\lexists} استفاده کنیم. اینجا باید به شرط متغیر ویژه توجه کنیم و
ثابتی را برگزینیم که نه در $\lexists[x][\lnot !A(x)]$ و نه در هیچ
فرضی که این عبارت به آن وابسته است رخ دهد. (اما چون هیچ !!{constant}sی
رخ نمی‌دهد، هر انتخابی مناسب است.)
\begin{prooftree}
\AxiomC{$\Discharge{\lexists[x][\lnot !A(x)]}{1}$}
\AxiomC{$\Discharge{\lnot !A(a)}{2}$}
\DeduceC{$\lnot \lforall[x][!A(x)]$}
\DischargeRule{\Elim{\lexists}}{2}
\BinaryInfC{$\lnot \lforall[x][!A(x)]$}
\DischargeRule{\Intro{\lif}}{1}
\UnaryInfC{$\lexists[x][\lnot !A(x)] \lif \lnot \lforall[x][!A(x)]$}
\end{prooftree}
برای اشتقاق $\lnot \lforall[x][!A(x)]$ می‌کوشیم از قاعدهٔ
\Intro{\lnot} استفاده کنیم: این کار مستلزم آن است که تناقضی را اشتقاق
کنیم، شاید با استفاده از $\lforall[x][!A(x)]$ به‌عنوان فرضی اضافی.
البته این تناقض ممکن است فرض $\lnot !A(a)$ را دربر گیرد که استنتاج
\Elim{\lexists} آن را رفع خواهد کرد. می‌توانیم آن را چنین ترتیب دهیم:
\begin{prooftree}
\AxiomC{$\Discharge{\lexists[x][\lnot !A(x)]}{1}$}
\AxiomC{$\Discharge{\lnot !A(a)}{2}, \Discharge{\lforall[x][!A(x)]}{3}$}
\DeduceC{$\lfalse$}
\DischargeRule{\Intro{\lnot}}{3}
\UnaryInfC{$\lnot \lforall[x][!A(x)]$}
\DischargeRule{\Elim{\lexists}}{2}
\BinaryInfC{$\lnot \lforall[x][!A(x)]$}
\DischargeRule{\Intro{\lif}}{1}
\UnaryInfC{$\lexists[x][\lnot !A(x)]\lif \lnot \lforall[x][!A(x)]$}
\end{prooftree}
به نظر می‌رسد به دستیابی به تناقض نزدیکیم. آسان‌ترین قاعده برای اعمال
\Elim{\lforall} است که هیچ شرط متغیر ویژه‌ای ندارد. چون می‌توانیم هر
ترمی را که می‌خواهیم به‌جای $x$ با سور همگانی به کار ببریم، معقول‌تر
است که همچنان از~$a$ استفاده کنیم تا به تناقض برسیم.
\begin{prooftree}
\AxiomC{$\Discharge{\lexists[x][\lnot !A(x)]}{1}$}
\AxiomC{$\Discharge{\lnot !A(a)}{2}$}
\AxiomC{$\Discharge{\lforall[x][!A(x)]}{3}$}
\RightLabel{\Elim{\lforall}}
\UnaryInfC{$!A(a)$}
\RightLabel{\Elim{\lnot}}
\BinaryInfC{$\lfalse$}
\DischargeRule{\Intro{\lnot}}{3}
\UnaryInfC{$\lnot \lforall[x][!A(x)]$}
\DischargeRule{\Elim{\lexists}}{2}
\BinaryInfC{$\lnot \lforall[x][!A(x)]$}
\DischargeRule{\Intro{\lif}}{1}
\UnaryInfC{$\lexists[x][\lnot !A(x)]\lif \lnot \lforall[x][!A(x)]$}
\end{prooftree}

به‌ویژه هنگام کار با سورها مهم است که در این مرحله دوباره بررسی کنیم
شرط متغیر ویژه نقض نشده باشد. چون تنها قاعده‌ای که اعمال کردیم و مشمول
شرط متغیر ویژه بود \Elim{\exists} است و متغیر ویژهٔ~$a$ در هیچ‌یک از
فرض‌هایی که به آن‌ها وابسته است رخ نمی‌دهد، این یک !!{derivation} درست
است.
\end{ex}

\begin{ex}
گاهی ممکن است !!a{formula} را از !!{formula}s دیگر اشتقاق کنیم. در این
موارد ممکن است فرض‌های رفع‌نشده داشته باشیم. مهم است که فرض‌هایمان و
نیز هدف نهایی را پیگیری کنیم.

ببینیم چگونه می‌توانیم !!a{derivation} از !!{formula}
$\lexists[x][!C(x,b)]$ از فرض‌های $\lexists[x][(!A(x)
\land !B(x))]$ و $\lforall[x][(!B(x) \lif !C(x,b))]$ به دست دهیم.
طبق معمول با نوشتن نتیجه در پایین آغاز می‌کنیم.
\begin{prooftree}
\AxiomC{}
\UnaryInfC{$\lexists[x][!C(x,b)]$}
\end{prooftree}

دو مقدمه برای کار داریم. برای استفاده از اولی، یعنی تلاش برای یافتن
!!a{derivation} از $\lexists[x][!C(x, b)]$ از $\lexists[x][(!A(x)
  \land !B(x))]$، باید از قاعدهٔ \Elim{\lexists} استفاده کنیم. چون این
قاعده شرط متغیر ویژه دارد، نخست آن را اعمال می‌کنیم. عبارت زیر را به
دست می‌آوریم:
\begin{prooftree}
\AxiomC{$\lexists[x][(!A(x) \land !B(x))]$}
\AxiomC{$\Discharge{!A(a) \land !B(a)}{1}$}
\DeduceC{$\lexists[x][!C(x,b)]$}
\DischargeRule{\Elim{\lexists}}{1}
\BinaryInfC{$\lexists[x][!C(x,b)]$}
\end{prooftree}
دو فرضی که با آن‌ها کار می‌کنیم در~$!B$ مشترک‌اند. شاید اکنون سودمند
باشد که \Elim{\land} را اعمال کنیم تا~$!B(a)$ را جدا سازیم.
\begin{prooftree}
\AxiomC{$\lexists[x][(!A(x) \land !B(x)])$}
\AxiomC{$\Discharge{!A(a)
\land !B(a)}{1}$}
\RightLabel{\Elim{\land}}
\UnaryInfC{$!B(a)$}
\DeduceC{$\lexists[x][!C(x,b)]$}
\DischargeRule{\Elim{\lexists}}{1}
\BinaryInfC{$\lexists[x][!C(x,b)]$}
\end{prooftree}

دومین فرضی که باید با آن کار کنیم~$\lforall[x][(!B(x) \lif
  !C(x,b))]$ است. چون شرط متغیر ویژه‌ای وجود ندارد، می‌توانیم با
استفاده از \Elim{\lforall}، $x$ را با !!{constant}~$a$ نمونه‌سازی کنیم
و $!B(a) \lif !C(a, b)$ را به دست آوریم. اکنون هم
$!B(a) \lif !C(a,b)$ و هم $!B(a)$ را داریم. حرکت بعدی ما باید کاربرد
سرراست قاعدهٔ \Elim{\lif} باشد.
\begin{prooftree}
\AxiomC{$\lexists[x][(!A(x) \land !B(x))]$}
\AxiomC{$\lforall[x][(!B(x) \lif !C(x,b))]$}
\RightLabel{\Elim{\lforall}}
\UnaryInfC{$!B(a) \lif !C(a,b)$}
\AxiomC{$\Discharge{!A(a)
\land !B(a)}{1}$}
\RightLabel{\Elim{\land}}
\UnaryInfC{$!B(a)$}
\RightLabel{\Elim{\lif}}
\BinaryInfC{$!C(a,b)$}
\DeduceC{$\lexists[x][!C(x,b)]$}
\DischargeRule{\Elim{\lexists}}{1}
\insertBetweenHyps{\hspace{-5em}}
\BinaryInfC{$\lexists[x][!C(x,b)]$}
\end{prooftree}
بسیار نزدیکیم!{} با یک کاربرد \Intro{\lexists} به هدفمان رسیده‌ایم.
\begin{prooftree}
\AxiomC{$\lexists[x][(!A(x) \land !B(x))]$}
\AxiomC{$\lforall[x][(!B(x) \lif !C(x,b))]$}
\RightLabel{\Elim{\lforall}}
\UnaryInfC{$!B(a) \lif !C(a,b)$}
\AxiomC{$\Discharge{!A(a)
\land !B(a)}{1}$}
\RightLabel{\Elim{\land}}
\UnaryInfC{$!B(a)$}
\RightLabel{\Elim{\lif}}
\BinaryInfC{$!C(a,b)$}
\RightLabel{\Intro{\lexists}}
\UnaryInfC{$\lexists[x][!C(x,b)]$}
\DischargeRule{\Elim{\lexists}}{1}
\insertBetweenHyps{\hspace{-5em}}
\BinaryInfC{$\lexists[x][!C(x,b)]$}
\end{prooftree}
چون در هر گام اطمینان حاصل کردیم که شرایط متغیر ویژه نقض نمی‌شوند،
می‌توانیم مطمئن باشیم این یک !!{derivation} درست است.
\end{ex}

\begin{ex}
!!a{derivation} از !!{formula}
$\lnot\lforall[x][!A(x)]$ از فرض‌های $\lforall[x][!A(x)]
\lif \lexists[y][!B(y)]$ و $\lnot\lexists[y][!B(y)]$ به دست دهید.
طبق معمول، !!{formula} هدف را در پایین می‌نویسیم.
\begin{prooftree}
\AxiomC{}
\UnaryInfC{$\lnot\lforall[x][!A(x)]$}
\end{prooftree}
آخرین خط !!{derivation} یک نقیض است، پس می‌کوشیم از \Intro{\lnot}
استفاده کنیم. برای این کار باید دریابیم چگونه تناقضی را !!{derive}
کنیم.
\begin{prooftree}
\AxiomC{$\Discharge{\lforall[x][!A(x)]}{1}$}
\DeduceC{$\lfalse$}
\DischargeRule{\Intro{\lnot}}{1}
\UnaryInfC{$\lnot\lforall[x][!A(x)]$}
\end{prooftree}
تا اینجا خوب پیش رفته‌ایم. می‌توانیم از \Elim{\lforall} استفاده کنیم،
اما آشکار نیست که این کار برای رسیدن به هدف کمکی کند. در عوض، از یکی از
فرض‌هایمان استفاده کنیم. $\lforall[x][!A(x)] \lif \lexists[y][!B(y)]$
همراه با $\lforall[x][!A(x)]$ به ما امکان می‌دهد قاعدهٔ \Elim{\lif}
را به کار ببریم.
\begin{prooftree}
\AxiomC{$\lforall[x][!A(x)] \lif \lexists[y][!B(y)]$}
\AxiomC{$\Discharge{\lforall[x][!A(x)]}{1}$}
\RightLabel{\Elim{\lif}}
\BinaryInfC{$\lexists[y][!B(y)]$}
\DeduceC{$\lfalse$}
\DischargeRule{\Intro{\lnot}}{1}
\UnaryInfC{$\lnot\lforall[x][!A(x)]$}
\end{prooftree}
اکنون یک فرض نهایی برای کار داریم و به نظر می‌رسد این فرض با استفاده
از \Elim{\lnot} ما را به تناقض می‌رساند.
\begin{prooftree}
\AxiomC{$\lnot\lexists[y][!B(y)]$}
\AxiomC{$\lforall[x][!A(x)] \lif \lexists[y][!B(y)]$}
\AxiomC{$\Discharge{\lforall[x][!A(x)]}{1}$}
\RightLabel{\Elim{\lif}}
\BinaryInfC{$\lexists[y][!B(y)]$}
\RightLabel{\Elim{\lnot}}
\BinaryInfC{$\lfalse$}
\DischargeRule{\Intro{\lnot}}{1}
\UnaryInfC{$\lnot\lforall[x][!A(x)]$}
\end{prooftree}
\end{ex}

\begin{prob}
!!{derivation}sی به دست دهید که موارد زیر را نشان دهند:
\begin{enumerate}
\item $\Proves (\lforall[x][!A(x)] \land \lforall[y][!B(y)]) \lif
\lforall[z][(!A(z) \land !B(z))]$.
\item $\Proves (\lexists[x][!A(x)] \lor \lexists[y][!B(y)]) \lif
\lexists[z][(!A(z) \lor !B(z))]$.
\item $\lforall[x][(!A(x) \lif !B)] \Proves \lexists[y][!A(y)] \lif !B$.
\item $\lforall[x][\lnot !A(x)] \Proves \lnot\lexists[x][!A(x)]$.
\item $\Proves \lnot\lexists[x][!A(x)] \lif \lforall[x][\lnot !A(x)]$.
\item $\Proves \lnot\lexists[x][\lforall[y][((!A(x,y) \lif \lnot
!A(y,y)) \land (\lnot !A(y,y) \lif !A(x,y)))]]$.
\end{enumerate}
\end{prob}

\begin{prob}
!!{derivation}sی به دست دهید که موارد زیر را نشان دهند:
\begin{enumerate}
\item $\Proves \lnot\lforall[x][!A(x)] \lif \lexists[x][\lnot!A(x)]$.
\item $(\lforall[x][!A(x)] \lif !B) \Proves \lexists[y][(!A(y) \lif !B)]$.
\item $\Proves \lexists[x][(!A(x) \lif \lforall[y][!A(y)])]$.
\end{enumerate}
(همهٔ این‌ها به قاعدهٔ~$\FalseCl$ نیاز دارند.)
\end{prob}

\end{document}
